\section{Infraestructura complementaria: un pipeline de MLOps para reentrenamiento continuo}

El capítulo de Alcance excluye explícitamente el aprendizaje continuo y el reentrenamiento automático embebido en el proceso de la aplicación anfitriona, por ser incompatible con el principio de operación autónoma sin dependencias externas que rige el diseño de logSguarDian en tiempo de ejecución. Esa exclusión sigue vigente: la librería npm entregada no reentrena modelos dentro del proceso de la aplicación que protege. Durante el desarrollo del proyecto, sin embargo, el equipo construyó y verificó, como infraestructura separada y exploratoria --no como parte del paquete npm ni de sus dependencias en tiempo de ejecución--, un pipeline de ocho fases para el ciclo completo de recolección de telemetría, reentrenamiento y despliegue controlado de nuevas versiones del modelo híbrido. Se documenta aquí como recomendación de arquitectura de producción para quien adopte logSguarDian a escala, y no como parte del alcance evaluado en los objetivos específicos de este trabajo.

La motivación de esta exploración es directa: cada cambio al pipeline de extracción de características exige un reentrenamiento completo del modelo, condición que el proyecto siguió de forma estricta a lo largo de sus más de diez ciclos de reentrenamiento documentados. Un despliegue real de logSguarDian a múltiples aplicaciones cliente enfrentaría el mismo problema de forma recurrente, y sin un mecanismo que sistematice la recolección de evidencia de producción, la curación de esa evidencia y la validación de un modelo candidato antes de promoverlo, cada ciclo de mejora tendría que repetirse manualmente. El pipeline construido --bajo el directorio \texttt{packages/mlops} del repositorio del proyecto-- explora si ese proceso puede automatizarse sin comprometer la garantía de bloqueo del modelo en producción.

Las ocho fases del pipeline son las siguientes:

\begin{enumerate}
    \item \textbf{Recolección de telemetría.} Un endpoint dedicado recibe únicamente el vector de 75 características de cada petición clasificada, nunca el payload crudo, siguiendo el mismo patrón de notificación asíncrona sin bloqueo (\textit{fire-and-forget}) ya empleado por el sistema de webhooks de la librería.
    \item \textbf{Simulación de flota.} Múltiples fuentes simuladas envían telemetría de forma concurrente para poblar un corpus de validación representativo de un despliegue con varios clientes.
    \item \textbf{Agrupamiento de anomalías.} El componente no supervisado (Isolation Forest) y un algoritmo de agrupamiento por densidad (DBSCAN) se aplican sobre el espacio de 63 características que consume el modelo no supervisado, para aislar patrones estadísticamente atípicos sin necesidad de una etiqueta previa.
    \item \textbf{Curación humana.} Una interfaz de línea de comandos presenta los grupos resultantes para revisión manual y genera un archivo con el mismo esquema que consume el script de unificación del corpus de entrenamiento del proyecto.
    \item \textbf{Reentrenamiento orquestado.} Un script coordina de punta a punta la unificación del corpus curado, el reentrenamiento del modelo híbrido y la generación de los artefactos ONNX correspondientes, reutilizando el mismo pipeline de entrenamiento documentado en el capítulo de Alcance.
    \item \textbf{Validación automática mediante cinco criterios de aceptación (\textit{gates}).} Se verifican, en orden creciente de costo computacional y con interrupción inmediata ante el primer fallo (\textit{fail-fast}): el contrato de features del modelo candidato, sus métricas de precisión y recall frente al conjunto de prueba bloqueado, la paridad de predicción entre el modelo en Python y su versión serializada en ONNX, su desempeño sobre la suite de pruebas end-to-end, y su huella de memoria bajo la arquitectura real de \textit{worker pool} de producción.
    \item \textbf{Despliegue canario.} Un modelo candidato que aprueba los cinco criterios anteriores se evalúa mediante dos mecanismos: reemplazo del corpus de prueba etiquetado como mecanismo primario de decisión, y una ventana acotada de sombra sobre tráfico en vivo como mecanismo secundario, nunca como proceso continuo indefinido.
    \item \textbf{Promoción o reversión.} La promoción de un modelo candidato a producción exige el cumplimiento simultáneo de un criterio estricto --cero casos en los que el candidato deje pasar un ataque que el modelo en producción bloquea correctamente-- y un criterio de concordancia general de al menos 99.5\% entre los veredictos de bloqueo del candidato y los del modelo vigente sobre el corpus de reemplazo.
\end{enumerate}

La verificación de este pipeline produjo un hallazgo concreto sobre el valor del componente no supervisado que merece documentarse de forma independiente. Sobre un corpus de validación de 244 eventos de telemetría simulados, el agrupamiento por Isolation Forest y DBSCAN aisló 17 grupos distintos. Uno de ellos, compuesto exclusivamente por seis peticiones enviadas desde una única fuente simulada dedicada a este propósito, obtuvo un puntaje de anomalía promedio de $-0.153$, muy por debajo del umbral de anomalía calibrado en producción ($0.0025$). Esas seis peticiones correspondían a un payload que combinaba deliberadamente, en una sola petición, sintaxis de las cuatro categorías de ataque cubiertas por logSguarDian --una cláusula SQL de eliminación de tabla, una etiqueta de script, una secuencia de ascenso de directorio y una sustitución de comando de shell--, un patrón nunca presente como tal en el corpus de entrenamiento, donde cada fuente enseña predominantemente una sola técnica. El componente supervisado clasificó las seis peticiones como una sola categoría de ataque (Cross-Site Scripting, la señal léxicamente más prominente del payload), sin indicio alguno de que la petición combinara además otras tres técnicas distintas. El agrupamiento no supervisado, en cambio, las aisló en su propio grupo, separado de los dieciséis grupos restantes y del grupo de ruido, sin haber recibido ninguna etiqueta previa que indicara que se trataba de una combinación de técnicas.

Este resultado no se obtuvo en el primer intento: una versión anterior del payload de prueba, que combinaba las mismas cuatro técnicas pero con menor grado de ofuscación simultánea, no produjo un puntaje de anomalía distinguible del resto del tráfico de ataque, y solo tras ajustar el payload para maximizar la combinación simultánea de longitud, entropía y coincidencias léxicas de las cuatro categorías a la vez se obtuvo la separación reportada. La conclusión correcta a partir de este hallazgo es acotada: el componente no supervisado detecta combinaciones de técnicas que se desvían lo suficiente de la distribución de entrenamiento en varios ejes simultáneamente, no cualquier técnica nueva de forma individual. Dentro de ese límite, el hallazgo confirma empíricamente el argumento presentado en el Marco Teórico sobre la complementariedad entre detección supervisada y no supervisada: el componente no supervisado identificó como atípica una combinación de ataques que el componente supervisado, por diseño, no tenía forma de distinguir de una categoría conocida.

Se recomienda, para quien adopte logSguarDian en un despliegue con múltiples aplicaciones cliente, evaluar la maduración de esta infraestructura complementaria --hoy una prueba de concepto verificada, no un servicio productivo-- como mecanismo de mantenimiento a largo plazo del modelo, manteniendo en todo momento la separación arquitectónica ya adoptada en este proyecto entre la librería embebida que protege la aplicación en tiempo real y la infraestructura externa que sostiene su ciclo de mejora continua.
